\centerline{\Huge \bf  Олимпиада II}
\centerline{\Huge \bf 6 класс}

\begin{flushright}
    \scriptsize{Заключительный отбор в Сириус 2022г.\\ Кто загуглит - тому бан!}
\end{flushright}

\problem{1} 
Можно ли квадрат $4 \times 4$ разрезать по клеточкам на четыре прямоугольника, среди которых нет равных? Напоминаем, что квадрат является частным случаем прямоугольника.


\problem{2}
Барон Мюнхгаузен утверждает, что он знает $2022$ числа, которые можно запи- сать в строчку так, что сумма любых двух соседних чисел будет отрицательна, а сумма любых трех подряд идущих – положительна. Можно ли ему верить? Ответ объясните.

\problem{3}
Каждую клетку таблицы $3 \times 2$ раскрашивают в один из трёх цветов так, чтобы соседние по стороне клетки были разных цветов. Сколько существует всевозмож- ных таблиц?

\problem{4}
Отрезок, длина которого $88$ см. разделён на четыре неравных отрезка. Расстояние между серединами самого левого и самого правого отрезков равно $46$ см. Найдите расстояние между серединами средних отрезков.

\problem{5}
У волшебника есть $5$ одинаковых закрытых шкатулок, расставленных по кругу, в одной из которых лежит бриллиант. У Пети есть чашечные весы без гирь, которые честно укажут, что шкатулка с бриллиантом тяжелее, если она есть на весах. За один ход Петя может взять любые две шкатулки, взвесить их на чашеч- ных весах. После взвешивания можно указать на любую шкатулку из пяти, если он уверен, что там есть бриллиант. Если же Петя не смог указать на шкатулку с бриллиантом, то волшебник произносит заклинание и бриллиант перемещается по часовой стрелке в соседнюю шкатулку. Если Петя ошибся, бриллиант исчезнет. Может ли Петя за два хода гарантированно найти шкатулку с бриллиантом?

\problem{6}
В квадрате $15 \times 15$ покрашено $222$ клетки. Требуется выбрать квадрат $n \times n$ в котором все клетки покрашены. Найдите наибольшее $n$ при котором можно выбрать указанный квадрат для любого расположения покрашенных клеток.
